\documentclass[../../main.tex]{subfiles}
\begin{document}
\section{常秩映射}\label{section:ch04-常秩映射}

线性映射最关键的线性代数性质是它的秩。事实上,正如\cref{theorem:appB-rank-canonical} 所示,如果我们可以对定义域和陪域独立地选取基,那么秩就是区分不同线性映射的唯一性质。

设 $M$ 与 $N$ 为带边或不带边的光滑流形。给定光滑映射 $F:M\to N$ 与点 $p\in M$,我们定义 $F$ 在 $p$ 处的秩为线性映射 $\mathrm{d}F_p:T_pM\to T_{F(p)}N$ 的秩;也就是 $F$ 在任一光滑坐标卡下的 Jacobian 矩阵的秩,或等价地为 $\operatorname{Im}\mathrm{d}F_p\subseteq T_{F(p)}N$ 的维数。如果 $F$ 在每一点都具有相同的秩 $r$,则称它具有常秩,并记作 $\operatorname{rank}F=r$。

由于线性映射的秩不会高于其定义域或陪域的维数(\cref{exercise:appB-rank-bound}),故 $F$ 在每一点处的秩都不超过 $\min\{\dim M,\dim N\}$。如果 $\mathrm{d}F_p$ 的秩等于这个上界,就称 $F$ 在 $p$ 处满秩;如果 $F$ 处处满秩,则称 $F$ 满秩。

最重要的常秩映射是满秩映射。若光滑映射 $F:M\to N$ 的微分在每一点都满射(等价地,$\operatorname{rank}F=\dim N$),则称它为\textbf{光滑淹没}。若其微分在每一点都单射(等价地,$\operatorname{rank}F=\dim M$),则称它为\textbf{光滑浸入}。

\begin{proposition}\label{proposition:ch04-full-rank-open}
设 $F:M\to N$ 是光滑映射, $p\in M$。若 $\mathrm{d}F_p$ 满射,则存在 $p$ 的邻域 $U$ 使得 $F|_U$ 是淹没。若 $\mathrm{d}F_p$ 单射,则存在 $p$ 的邻域 $U$ 使得 $F|_U$ 是浸入。
\end{proposition}

\begin{proof}
若我们在 $p$ 附近为 $M$、在 $F(p)$ 附近为 $N$ 任选光滑坐标,则任一假设都意味着 $F$ 在坐标下的 Jacobian 矩阵在 $p$ 处满秩。\cref{example:ch01-open-full-rank-matrices} 表明满秩的 $m\times n$ 矩阵之集是 $M(m\times n,\mathbb{R})$ 的开子集(其中 $m=\dim M$, $n=\dim N$),故由连续性, $F$ 的 Jacobian 在 $p$ 的某邻域内均满秩。
\end{proof}

正如本章将看到的,光滑淹没与光滑浸入分别局部地表现得像满射线性映射与单射线性映射。(还存在拓扑淹没与拓扑浸入的类似概念,它们适用于仅连续的映射。我们并不需要使用这些概念,但为完整起见,将在本章后面加以描述。)

\begin{example}[Submersions and Immersions]\label{example:ch04-submersions-immersions}
设 $M_1,\dots,M_k$ 为光滑流形。则每个投影映射 $\pi_i:M_1\times\cdots\times M_k\to M_i$ 都是光滑淹没。特别地,投影 $\pi:\mathbb{R}^{n+k}\to\mathbb{R}^n$(投向前 $n$ 个坐标)是光滑淹没。

若 $\gamma:J\to M$ 是带边或不带边光滑流形 $M$ 中的光滑曲线,则 $\gamma$ 是光滑浸入当且仅当对所有 $t\in J$ 都有 $\gamma'(t)\neq 0$。

若 $M$ 是光滑流形,且其切丛 $TM$ 被赋予\cref{proposition:ch03-tangent-bundle-structure} 所述的光滑流形结构,则投影 $\pi:TM\to M$ 是光滑淹没。为验证这一点,只需注意:关于 $M$ 的开子集 $U$ 上的任一光滑局部坐标 $\bigl(x^i\bigr)$ 以及 $\pi^{-1}(U)\subseteq TM$ 上相应的自然坐标 $\bigl(x^i,v^i\bigr)$(见\cref{proposition:ch03-tangent-bundle-structure}), $\pi$ 的坐标表示为 $\widehat{\pi}(x,v)=x$。

由下式给出的光滑映射 $X:\mathbb{R}^2\to\mathbb{R}^3$
\begin{align*}
X(u,v)=\bigl((2+\cos 2\pi u)\cos 2\pi v,\,(2+\cos 2\pi u)\sin 2\pi v,\,\sin 2\pi u\bigr)
\end{align*}
是 $\mathbb{R}^2$ 到 $\mathbb{R}^3$ 的光滑浸入,其像是由 $(y,z)$ 平面上的圆 $(y-2)^2+z^2=1$ 绕 $z$ 轴旋转所得的甜甜圈状曲面(\cref{figure:ch04-torus-revolution})。
\end{example}

\begin{figure}[ht]
\centering
\includegraphics[width=0.8\textwidth]{fig-04-1.png}
\caption{$\mathbb{R}^3$ 中的旋转环面}\label{figure:ch04-torus-revolution}
\end{figure}

\begin{exercise}\label{exercise:ch04-verify-42}
验证上一例中所给出的断言。
\end{exercise}

\begin{exercise}\label{exercise:ch04-composition-rank}
证明光滑淹没的复合仍是光滑淹没,光滑浸入的复合仍是光滑浸入。举一个反例说明:常秩映射的复合未必具有常秩。
\end{exercise}

\subsection{局部微分同胚}\label{subsection:ch04-local-diffeomorphisms}

若 $M$ 与 $N$ 为带边或不带边的光滑流形,称映射 $F:M\to N$ 为局部微分同胚,如果每个点 $p\in M$ 都有邻域 $U$,使得 $F(U)$ 在 $N$ 中开,且 $F|_U:U\to F(U)$ 是微分同胚。下一定理是局部微分同胚最重要性质的关键。

\begin{theorem}[Inverse Function Theorem for Manifolds]\label{theorem:ch04-inverse-function}
设 $M$ 与 $N$ 为光滑流形, $F:M\to N$ 为光滑映射。若 $p\in M$ 是使 $\mathrm{d}F_p$ 可逆的点,则存在 $p$ 的连通邻域 $U_0$ 与 $F(p)$ 的连通邻域 $V_0$,使得 $F|_{U_0}:U_0\to V_0$ 是微分同胚。
\end{theorem}

\begin{proof}
$\mathrm{d}F_p$ 双射这一事实意味着 $M$ 与 $N$ 维数相同,设为 $n$。选取以 $p$ 为中心的光滑坐标卡 $(U,\varphi)$ 与以 $F(p)$ 为中心的光滑坐标卡 $(V,\psi)$,并使 $F(U)\subseteq V$。于是 $\widehat{F}=\psi\circ F\circ\varphi^{-1}$ 是从开子集 $\widehat{U}=\varphi(U)\subseteq\mathbb{R}^n$ 到 $\widehat{V}=\psi(V)\subseteq\mathbb{R}^n$ 的光滑映射,且 $\widehat{F}(0)=0$。由于 $\varphi$ 与 $\psi$ 是微分同胚,微分
\begin{align*}
\mathrm{d}\widehat{F}_0=\mathrm{d}\psi_{F(p)}\circ\mathrm{d}F_p\circ\mathrm{d}\bigl(\varphi^{-1}\bigr)_0
\end{align*}
非奇异。通常的反函数定理(\cref{theorem:appC-inverse-function})表明:存在含 $0$ 的连通开子集 $\widehat{U}_0\subseteq\widehat{U}$ 与 $\widehat{V}_0\subseteq\widehat{V}$,使得 $\widehat{F}$ 限制为从 $\widehat{U}_0$ 到 $\widehat{V}_0$ 的微分同胚。于是 $U_0=\varphi^{-1}(\widehat{U}_0)$ 与 $V_0=\psi^{-1}(\widehat{V}_0)$ 分别是 $p$ 与 $F(p)$ 的连通邻域,由复合即知 $F|_{U_0}$ 是从 $U_0$ 到 $V_0$ 的微分同胚。
\end{proof}

需要注意,我们只对不带边流形陈述了\cref{theorem:ch04-inverse-function}。事实上,当定义域具有非空边界时,它可能失效(见习题 4-1)。然而,当陪域具有非空边界时,仍有有用的结论:只要映射取值于陪域的内部,同样的结论成立,因为内部是不带边的光滑流形。习题 4-2 表明:在微分可逆的点处,情况总是如此。

\begin{proposition}[Elementary Properties of Local Diffeomorphisms]\label{proposition:ch04-local-diffeo-properties}
\begin{enumerate}[(a)]
\item 局部微分同胚的复合仍是局部微分同胚。
\item 光滑流形之间有限个局部微分同胚的积仍是局部微分同胚。
\item 每个局部微分同胚都是局部同胚与开映射。
\item 局部微分同胚到带边或不带边开子流形上的限制仍是局部微分同胚。
\item 每个微分同胚都是局部微分同胚。
\item 每个双射局部微分同胚都是微分同胚。
\item 带边或不带边光滑流形之间的映射是局部微分同胚,当且仅当在其定义域每一点附近,它都具有一个为局部微分同胚的坐标表示。
\end{enumerate}
\end{proposition}

\begin{exercise}\label{exercise:ch04-prove-local-diffeo-props}
证明上述命题。
\end{exercise}

\begin{proposition}\label{proposition:ch04-local-diffeo-equiv}
设 $M$ 与 $N$ 为(不带边)光滑流形, $F:M\to N$ 为映射。
\begin{enumerate}[(a)]
\item $F$ 是局部微分同胚,当且仅当它同时是光滑浸入与光滑淹没。
\item 若 $\dim M=\dim N$,且 $F$ 是光滑浸入或光滑淹没之一,则它是局部微分同胚。
\end{enumerate}
\end{proposition}

\begin{proof}
先设 $F$ 是局部微分同胚。给定 $p\in M$,存在 $p$ 的邻域 $U$ 使得 $F$ 是从 $U$ 到 $F(U)$ 的微分同胚。于是由\cref{proposition:ch03-diffeomorphism-differential}(d) 知 $\mathrm{d}F_p:T_pM\to T_{F(p)}N$ 是同构。因此 $\operatorname{rank}F=\dim M=\dim N$,故 $F$ 同时是光滑浸入与光滑淹没。反之,若 $F$ 同时是光滑浸入与光滑淹没,则 $\mathrm{d}F_p$ 在每个 $p\in M$ 处都是同胚,而流形的反函数定理(\cref{theorem:ch04-inverse-function})表明 $p$ 有一邻域,使 $F$ 在其上限制为到其像上的微分同胚。这就证明了 (a)。

为证 (b),注意到当 $M$ 与 $N$ 维数相同时, $\mathrm{d}F_p$ 的单射性或满射性都蕴涵双射性,故 $F$ 是光滑淹没当且仅当它是光滑浸入,从而 (b) 由 (a) 推出。
\end{proof}

\begin{exercise}\label{exercise:ch04-local-diffeo-boundary}
证明:若允许 $N$ 为带边光滑流形,则\cref{proposition:ch04-local-diffeo-equiv} 的结论仍成立;但若允许 $M$ 为带边流形,则不成立。(见习题 4-1 与习题 4-2。)
\end{exercise}

\begin{exercise}\label{exercise:ch04-local-diffeo-universal}
设 $M,N,P$ 为带边或不带边的光滑流形, $F:M\to N$ 是局部微分同胚。证明下列结论:
\begin{enumerate}[(a)]
\item 若 $G:P\to M$ 连续,则 $G$ 光滑当且仅当 $F\circ G$ 光滑。
\item 若进一步 $F$ 满射,且 $G:N\to P$ 为任一映射,则 $G$ 光滑当且仅当 $G\circ F$ 光滑。
\end{enumerate}
\end{exercise}

\begin{example}[Local Diffeomorphisms]\label{example:ch04-local-diffeomorphisms}
\cref{example:ch01-circle-covering}(b) 中定义的映射 $\varepsilon:\mathbb{R}\to S^1$ 是局部微分同胚,因为在每一点附近它都有形如 $t\mapsto 2\pi t+c$ 的坐标表示,而这是局部微分同胚。类似地, \cref{example:ch01-circle-covering}(c) 中定义的映射 $\varepsilon^n:\mathbb{R}^n\to\mathbb{T}^n$ 是局部微分同胚,因为它是若干局部微分同胚的积。
\end{example}

在本章末尾,我们将考察一类重要的特殊局部微分同胚,即光滑覆叠映射。

\subsection{秩定理}\label{subsection:ch04-rank-theorem}

关于常秩映射最重要的事实是反函数定理的下述推论,它说明:经过坐标变换,常秩光滑映射可被局部化成一种特别简单的典范形式。它是\cref{theorem:appB-rank-canonical} 所给出的线性映射典范形式定理的非线性版本。

\begin{theorem}[Rank Theorem]\label{theorem:ch04-rank-theorem}
设 $M$ 与 $N$ 分别为 $m$ 维与 $n$ 维光滑流形, $F:M\to N$ 是秩为常数 $r$ 的光滑映射。对每个 $p\in M$,存在以 $p$ 为中心的 $M$ 的光滑坐标卡 $(U,\varphi)$ 与以 $F(p)$ 为中心的 $N$ 的光滑坐标卡 $(V,\psi)$,使 $F(U)\subseteq V$,且在此坐标下 $F$ 具有如下坐标表示
\begin{align}\label{align:ch04-rank-canonical}
\widehat{F}\bigl(x^1,\dots,x^r,x^{r+1},\dots,x^m\bigr)=\bigl(x^1,\dots,x^r,0,\dots,0\bigr).
\end{align}
特别地,若 $F$ 是光滑淹没,则上式成为
\begin{align}\label{align:ch04-rank-submersion}
\widehat{F}\bigl(x^1,\dots,x^n,x^{n+1},\dots,x^m\bigr)=\bigl(x^1,\dots,x^n\bigr),
\end{align}
而若 $F$ 是光滑浸入,则为
\begin{align}\label{align:ch04-rank-immersion}
\widehat{F}\bigl(x^1,\dots,x^m\bigr)=\bigl(x^1,\dots,x^m,0,\dots,0\bigr).
\end{align}
\end{theorem}

\begin{proof}
由于定理是局部的,在选取光滑坐标后,我们可将 $M$ 与 $N$ 分别替换为开子集 $U\subseteq\mathbb{R}^m$ 与 $V\subseteq\mathbb{R}^n$。$\mathrm{D}F(p)$ 的秩为 $r$ 这一事实意味着其矩阵有某个 $r\times r$ 子行列式非零。通过重排坐标,可设它是左上子矩阵 $\bigl(\partial F^i/\partial x^j\bigr)$,其中 $i,j=1,\dots,r$。把 $\mathbb{R}^m$ 中的标准坐标重记为 $(x,y)=\bigl(x^1,\dots,x^r,y^1,\dots,y^{m-r}\bigr)$,把 $\mathbb{R}^n$ 中的重记为 $(v,w)=\bigl(v^1,\dots,v^r,w^1,\dots,w^{n-r}\bigr)$。

通过坐标的初始平移,可不失一般性地设 $p=(0,0)$ 且 $F(p)=(0,0)$。若把 $F(x,y)$ 写为 $\bigl(Q(x,y),R(x,y)\bigr)$,其中 $Q:U\to\mathbb{R}^r$ 与 $R:U\to\mathbb{R}^{n-r}$ 为光滑映射,则我们的假设是 $\bigl(\partial Q^i/\partial x^j\bigr)$ 在 $(0,0)$ 处非奇异。

定义 $\varphi:U\to\mathbb{R}^m$ 为 $\varphi(x,y)=\bigl(Q(x,y),y\bigr)$。它在 $(0,0)$ 处的全导数为
\begin{align*}
\mathrm{D}\varphi(0,0)=
\begin{pmatrix}
\dfrac{\partial Q^i}{\partial x^j}(0,0) & \dfrac{\partial Q^i}{\partial y^j}(0,0)\\
0 & \delta^i_j
\end{pmatrix},
\end{align*}
此处使用如下标准记号:对正整数 $i$ 与 $j$,符号 $\delta^i_j$ 称为 Kronecker 符号,定义为
\begin{align}\label{align:ch04-kronecker}
\delta^i_j=
\begin{cases}
1,& i=j,\\
0,& i\neq j.
\end{cases}
\end{align}
由假设,矩阵 $\mathrm{D}\varphi(0,0)$ 非奇异。因此由反函数定理,存在 $(0,0)$ 的连通邻域 $U_0$ 与 $\varphi(0,0)=(0,0)$ 的连通邻域 $\widetilde{U}_0$,使得 $\varphi:U_0\to\widetilde{U}_0$ 是微分同胚。必要时缩小 $U_0$ 与 $\widetilde{U}_0$,可设 $\widetilde{U}_0$ 是开立方体。把逆映射写成 $\varphi^{-1}(x,y)=\bigl(A(x,y),B(x,y)\bigr)$,其中 $A:\widetilde{U}_0\to\mathbb{R}^r$ 与 $B:\widetilde{U}_0\to\mathbb{R}^{m-r}$ 为光滑函数,我们计算
\begin{align}\label{align:ch04-phi-inv}
(x,y)=\varphi\bigl(A(x,y),B(x,y)\bigr)=\bigl(Q\bigl(A(x,y),B(x,y)\bigr),B(x,y)\bigr).
\end{align}
比较 $y$ 分量得 $B(x,y)=y$,因此 $\varphi^{-1}$ 具有形式
\begin{align*}
\varphi^{-1}(x,y)=\bigl(A(x,y),y\bigr).
\end{align*}
另一方面, $\varphi\circ\varphi^{-1}=\operatorname{Id}$ 蕴涵 $Q\bigl(A(x,y),y\bigr)=x$,从而 $F\circ\varphi^{-1}$ 具有形式
\begin{align*}
F\circ\varphi^{-1}(x,y)=\bigl(x,\widetilde{R}(x,y)\bigr),
\end{align*}
其中 $\widetilde{R}:\widetilde{U}_0\to\mathbb{R}^{n-r}$ 由 $\widetilde{R}(x,y)=R\bigl(A(x,y),y\bigr)$ 定义。这个复合映射在任一点 $(x,y)\in\widetilde{U}_0$ 处的 Jacobian 矩阵为
\begin{align*}
\mathrm{D}\bigl(F\circ\varphi^{-1}\bigr)(x,y)=
\begin{pmatrix}
\delta^i_j & 0\\
\dfrac{\partial \widetilde{R}^i}{\partial x^j}(x,y) & \dfrac{\partial \widetilde{R}^i}{\partial y^j}(x,y)
\end{pmatrix}.
\end{align*}
由于与微分同胚复合不改变映射的秩,该矩阵在 $\widetilde{U}_0$ 中处处秩为 $r$。前 $r$ 列显然线性无关,故秩能为 $r$ 的必要条件是导数 $\partial\widetilde{R}^i/\partial y^j$ 在 $\widetilde{U}_0$ 上恒为零,这蕴涵 $\widetilde{R}$ 实际上不依赖于 $\bigl(y^1,\dots,y^{m-r}\bigr)$。(这正是我们把 $\widetilde{U}_0$ 取成立方体的原因之一。)于是若令 $S(x)=\widetilde{R}(x,0)$,则有
\begin{align}\label{align:ch04-after-phi}
F\circ\varphi^{-1}(x,y)=\bigl(x,S(x)\bigr).
\end{align}

为完成证明,我们需在 $V$ 中 $(0,0)$ 的某邻域内定义一个合适的光滑坐标卡。设 $V_0\subseteq V$ 为开子集
\begin{align*}
V_0=\{(v,w)\in V:(v,0)\in\widetilde{U}_0\}.
\end{align*}
则 $V_0$ 是 $(0,0)$ 的邻域。由于 $\widetilde{U}_0$ 是立方体且 $F\circ\varphi^{-1}$ 具有形式\eqref{align:ch04-after-phi},可知 $F\circ\varphi^{-1}(\widetilde{U}_0)\subseteq V_0$,从而 $F(U_0)\subseteq V_0$。定义 $\psi:V_0\to\mathbb{R}^n$ 为 $\psi(v,w)=\bigl(v,w-S(v)\bigr)$。这是到其像上的微分同胚,因为其逆由 $\psi^{-1}(s,t)=\bigl(s,t+S(s)\bigr)$ 显式给出;因而 $(V_0,\psi)$ 是光滑坐标卡。由\eqref{align:ch04-after-phi} 得
\begin{align*}
\psi\circ F\circ\varphi^{-1}(x,y)=\psi\bigl(x,S(x)\bigr)=\bigl(x,S(x)-S(x)\bigr)=(x,0),
\end{align*}
此即所证。
\end{proof}

下一个推论可视为秩定理的一个更具不变性的陈述。它说:常秩映射恰好就是那些局部行为与其微分相同的映射。

\begin{corollary}\label{corollary:ch04-constant-rank-linear}
设 $M$ 与 $N$ 为光滑流形, $F:M\to N$ 为光滑映射,并设 $M$ 连通。则下列条件等价:
\begin{enumerate}[(a)]
\item 对每个 $p\in M$,存在包含 $p$ 与 $F(p)$ 的光滑坐标卡,使 $F$ 的坐标表示为线性映射。
\item $F$ 具有常秩。
\end{enumerate}
\end{corollary}

\begin{proof}
先设 $F$ 在每一点附近都有线性坐标表示。由于每个线性映射都有常秩,故 $F$ 的秩在每点附近为常数,从而由连通性知它在整个 $M$ 上为常数。反之,若 $F$ 具有常秩,秩定理表明它在每点附近都有线性坐标表示\eqref{align:ch04-rank-canonical}。
\end{proof}

秩定理是纯局部的陈述。然而,它有如下强有力的整体推论。

\begin{theorem}[Global Rank Theorem]\label{theorem:ch04-global-rank}
设 $M$ 与 $N$ 为光滑流形,并设 $F:M\to N$ 是常秩光滑映射。
\begin{enumerate}[(a)]
\item 若 $F$ 满射,则它是光滑淹没。
\item 若 $F$ 单射,则它是光滑浸入。
\item 若 $F$ 双射,则它是微分同胚。
\end{enumerate}
\end{theorem}

\begin{proof}
令 $m=\dim M$, $n=\dim N$,并设 $F$ 具有常秩 $r$。为证 (a),反设 $F$ 不是光滑淹没,即 $r<n$。由秩定理,对每个 $p\in M$,存在以 $p$ 为中心的 $M$ 的坐标卡 $(U,\varphi)$ 与以 $F(p)$ 为中心的 $N$ 的坐标卡 $(V,\psi)$,使 $F(U)\subseteq V$ 且 $F$ 的坐标表示由\eqref{align:ch04-rank-canonical} 给出。(见\cref{figure:ch04-global-rank}。)必要时缩小 $U$,可设它是正则坐标球且 $F(\overline{U})\subseteq V$。这蕴涵 $F(\overline{U})$ 是集合 $\{y\in V:y^{r+1}=\cdots=y^n=0\}$ 的紧子集,故它在 $N$ 中闭且不含 $N$ 的任何开子集;因而它在 $N$ 中无处稠密。由于流形的每个开覆盖都有可数子覆盖,我们可选出可数多个这样的坐标卡 $\{(U_i,\varphi_i)\}$ 覆盖 $M$,并有相应坐标卡 $\{(V_i,\psi_i)\}$ 覆盖 $F(M)$。由于 $F(M)$ 等于无处稠密集合 $F(\overline{U_i})$ 的可数并,由 Baire 纲定理(\cref{theorem:appA-baire})知 $F(M)$ 在 $N$ 中内部为空,这意味着 $F$ 不可能满射。

为证 (b),反设 $F$ 不是光滑浸入,即 $r<m$。由秩定理,对每个 $p\in M$,我们可在 $p$ 与 $F(p)$ 附近选取坐标卡,使 $F$ 具有坐标表示\eqref{align:ch04-rank-canonical}。于是对任意充分小的 $\varepsilon$,都有 $F(0,\dots,0,\varepsilon)=F(0,\dots,0,0)$,故 $F$ 不单射。

最后, (c) 由 (a) 与 (b) 推出,因为双射常秩光滑映射由 (a) 是光滑淹没、由 (b) 是光滑浸入;于是\cref{proposition:ch04-local-diffeo-equiv} 表明 $F$ 是局部微分同胚,又因它双射,故它是微分同胚。
\end{proof}

\begin{figure}[ht]
\centering
\includegraphics[width=0.8\textwidth]{fig-04-2.png}
\caption{定理 4.14(a) 的证明}\label{figure:ch04-global-rank}
\end{figure}

\subsection{带边流形的秩定理}\label{subsection:ch04-rank-boundary}

在带边流形的背景下,我们只在一种特殊情形下需要秩定理:即定义域为带边光滑流形的光滑浸入。当然,由于带边光滑流形的内部本身是光滑流形,在定义域的任一内点附近,通常的秩定理都适用。对于边界点,我们有如下的秩定理替代。

\begin{theorem}[Local Immersion Theorem for Manifolds with Boundary]\label{theorem:ch04-immersion-boundary}
设 $M$ 是带边光滑 $m$ 流形, $N$ 是光滑 $n$ 流形, $F:M\to N$ 是光滑浸入。对任一 $p\in\partial M$,存在以 $p$ 为中心的 $M$ 的光滑边界坐标卡 $(U,\varphi)$ 与以 $F(p)$ 为中心的 $N$ 的光滑坐标卡 $(V,\psi)$,使 $F(U)\subseteq V$,且在此坐标下 $F$ 具有坐标表示
\begin{align}\label{align:ch04-immersion-boundary}
\widehat{F}\bigl(x^1,\dots,x^m\bigr)=\bigl(x^1,\dots,x^m,0,\dots,0\bigr).
\end{align}
\end{theorem}

\begin{proof}
通过先为 $M$ 与 $N$ 选取光滑坐标卡,可设 $M$ 与 $N$ 分别是 $\mathbb{H}^m$ 与 $\mathbb{R}^n$ 的开子集,并设 $p=0\in\mathbb{H}^m$ 且 $F(p)=0\in\mathbb{R}^n$。由 $\mathbb{H}^m$ 上函数光滑性的定义, $F$ 可延拓为光滑映射 $\widetilde{F}:W\to\mathbb{R}^n$,其中 $W$ 是 $\mathbb{R}^m$ 中含 $0$ 的某开子集。由于 $\mathrm{d}\widetilde{F}_0=\mathrm{d}F_0$ 单射,必要时缩小 $W$,可设 $\widetilde{F}$ 是光滑浸入。把 $\mathbb{R}^m$ 上的坐标记为 $x=\bigl(x^1,\dots,x^m\bigr)$,把 $\mathbb{R}^n$ 上的坐标记为 $(v,w)=\bigl(v^1,\dots,v^m,w^1,\dots,w^{n-m}\bigr)$。

由秩定理,存在以 $0$ 为中心的 $\mathbb{R}^m$ 的光滑坐标卡 $(U_0,\varphi_0)$ 与以 $0$ 为中心的 $\mathbb{R}^n$ 的光滑坐标卡 $(V_0,\psi_0)$,使 $\widehat{F}=\psi_0\circ\widetilde{F}\circ\varphi_0^{-1}$ 由\eqref{align:ch04-immersion-boundary} 给出。这些坐标唯一的问题是 $\varphi_0$ 未必限制为 $M$ 的边界坐标卡。但我们可如下容易地修正它。由于 $\varphi_0$ 是从 $U_0$ 到开子集 $\widehat{U}_0=\varphi_0(U_0)\subseteq\mathbb{R}^m$ 的微分同胚,映射 $\varphi_0^{-1}\times\operatorname{Id}_{\mathbb{R}^{n-m}}$ 是从 $\widehat{U}_0\times\mathbb{R}^{n-m}$ 到 $U_0\times\mathbb{R}^{n-m}$ 的微分同胚。令
\begin{align*}
\xi=\bigl(\varphi_0^{-1}\times\operatorname{Id}_{\mathbb{R}^{n-m}}\bigr)\circ\psi_0,
\end{align*}
它是从 $V_0$ 中含 $0$ 的某开子集 $V$ 到 $\mathbb{R}^n$ 中 $0$ 的某邻域的微分同胚。利用\eqref{align:ch04-immersion-boundary},我们计算
\begin{align*}
\xi\circ F(x)
&=\bigl(\varphi_0^{-1}\times\operatorname{Id}_{\mathbb{R}^{n-m}}\bigr)\circ\psi_0\circ F\circ\varphi_0^{-1}\circ\varphi_0(x)\\
&=\bigl(\varphi_0^{-1}\times\operatorname{Id}_{\mathbb{R}^{n-m}}\bigr)\circ\widehat{F}\bigl(\varphi_0(x)\bigr)\\
&=\bigl(\varphi_0^{-1}\times\operatorname{Id}_{\mathbb{R}^{n-m}}\bigr)\bigl(\varphi_0(x),0\bigr)=(x,0).
\end{align*}
因此, $M$ 的原始坐标(限制到 $0$ 的充分小邻域)与 $N$ 的坐标卡 $(V,\xi)$ 即满足所需条件。
\end{proof}

对于从带边流形出发的更一般常秩映射,也可以证明类似的定理,但其证明更为繁琐,因为 $F$ 到开子集的延拓并不自动具有常秩。由于我们不需要这一更一般的结果,留待有兴趣的读者自行探讨(习题 4-3)。另一方面,当陪域是带边流形时,情况要复杂得多:由于映射的像可能以无法预料的方式与边界相交,若不对映射施加强限制,就无法把它化成简单的典范形式。

\end{document}